\section{Elementary structural results}

\subsection{Odd orders and group tables}

\begin{proposition}\label{prop:odd-order}
Every Latin square of odd order $n>1$ has pattern $\TTT$.
\end{proposition}

\begin{proof}
By Lemma~\ref{lem:fixed-point-free}, every induced permutation between two
distinct lines is fixed-point-free on an odd set.  If all of its nontrivial
cycles were even, their lengths would have even sum.  Hence it has an odd
cycle of length greater than one.  This argument applies separately to any
pair of distinct rows, columns, and symbols.
\end{proof}

Let $G$ be a finite group and let $L_G(x,y)=xy$ be its Cayley table.

\begin{proposition}\label{prop:group-tables}
The row permutation induced by rows $a,b$ of $L_G$ is left multiplication by
$b^{-1}a$.  The column permutation induced by columns $c,d$ is right
multiplication by $cd^{-1}$.  The symbol permutation induced by symbols
$u,v$ is right multiplication by $u^{-1}v$.  Each cycle length is therefore
the order of the corresponding group element.
\end{proposition}

\begin{proof}
The row identity is
$r_b^{-1}r_a(x)=b^{-1}ax$.  Likewise
$q_d^{-1}q_c(x)=xcd^{-1}$.  Finally, $s_u(c)=uc^{-1}$ and
$s_v^{-1}(a)=a^{-1}v$, whence
$s_v^{-1}s_u(c)=cu^{-1}v$.  Left or right multiplication by $g$ has orbits
equal to cosets of $\langle g\rangle$, all of length $\operatorname{ord}(g)$.
\end{proof}

\begin{corollary}\label{cor:group-fff}
The Cayley table of a finite group $G$ is FFF if and only if $G$ is a
$2$-group.
\end{corollary}

\begin{proof}
If $G$ is a $2$-group, every nonidentity element has even order, so
Proposition~\ref{prop:group-tables} gives no nontrivial odd cycle.  Conversely,
if $|G|$ has an odd prime divisor, Cauchy's theorem supplies a nonidentity
element of odd order; the corresponding induced permutations give a failing
view.
\end{proof}

\subsection{A closure criterion for group isotopy}

Fix a base row $b$ of a Latin square $L$ and define
\[
 H_L(b)=\{r_b^{-1}r_x:x\in X\}\subseteq\operatorname{Sym}(X).
\]

\begin{proposition}[Basis-family closure criterion]\label{prop:closure}
A Latin square $L$ is isotopic to a group table if and only if $H_L(b)$ is
closed under composition.  This condition is independent of the chosen base
row.
\end{proposition}

\begin{proof}
Suppose first that $H=H_L(b)$ is closed.  It contains the identity, has
$|X|$ distinct elements, and is a finite submonoid of a symmetric group;
hence it is a group.  For each $z\in X$, the values
$(r_b^{-1}r_x)(z)$, as $x$ varies, are all distinct, since a column of $L$
contains every symbol once.  Thus $H$ acts regularly on $X$.

Fix $z_0\in X$.  Relabel rows by $x\mapsto h_x=r_b^{-1}r_x$.  Use the regular
action to identify every column $z$ with the unique $h_z\in H$ satisfying
$h_z(z_0)=z$, and identify every relabeled symbol $r_b^{-1}(u)$ in the same
way.  The entry in row $h_x$ and column $h_z$ is represented by the unique
element sending $z_0$ to $h_xh_z(z_0)$, namely $h_xh_z$.  The resulting
isotope is therefore the Cayley table of $H$.

Conversely, for a group table the family at base row $b$ is the full left
regular representation $\{\lambda_{b^{-1}x}:x\in G\}$ and is closed.
Isotopy conjugates the family and relabels its members, preserving closure.
The same argument for any base row proves base independence.
\end{proof}

The same criterion can equivalently be written in the column or symbol view,
by parastrophy.
